%%%%%%%%%%%%%%%%%%%%%%%%%%%%%%%%%%%%%%%%%%%%%%%%%%%
%
%  CHAPTER V -- COMPLETED RESEARCH III:
%  AN ALGEBRAIC FRAMEWORK FOR VOLUMETRIC NONWOVEN
%  TEXTILES AND PERIODIC PATTERN DESIGN
%
%%%%%%%%%%%%%%%%%%%%%%%%%%%%%%%%%%%%%%%%%%%%%%%%%%%

\chapter{\uppercase{Completed Research III: An Algebraic Framework for Volumetric Nonwoven Textiles and Periodic Pattern Design}}
\label{cha:nonwoven}

Chapter~\ref{cha:volumetric-bravais} treated volumetric weaves in which threads strictly follow an over--under interlacing rule at every cell interface. This chapter summarizes~\cite{yildiz2024nonwoven}, which relaxes that rule and, in doing so, both (i) extends the cell-induced framework to \emph{nonwoven} structures and (ii) introduces the algebraic machinery -- a numeral system based on $2^K$-adic and $8^K$-adic ``zeros'' -- that this dissertation uses to \emph{design} within a periodic configuration space directly, rather than enumerating it after the fact.

\section{Motivation: Beyond the Over--Under Rule}

Woven fabrics are only one way of achieving fabric-like connectivity. \emph{Nonwoven} structures -- more precisely, planar nonwoven textiles that mimic the connectivity of woven fabric without the crossing rule -- instead form a network of paths joined by local attachments (nails, gluing, welding, or chemical bonding) rather than physical interlacing. Extending nonwoven connectivity into three dimensions turns out to be comparatively direct, in contrast to the significant conceptual step required to extend \emph{woven} structures into three dimensions (Chapter~\ref{cha:volumetric-bravais}). This asymmetry is itself useful: volumetric nonwoven structures can be designed to mimic the mechanical behavior of woven fabrics while also realizing configurations -- including representations of previously unknown infinite regular and semi-regular polyhedra, and flexible ``wood''-like frame structures -- that are simply not achievable with planar or volumetric woven constructions.

\section{Representation: From Matrices to Octal Tensors}

Two-way, two-layer nonwoven structures are represented with a binary matrix recording, for each grid position, whether a connection is present between two layers -- a direct generalization of Gr\"unbaum and Shephard's matrix notation for planar weaves~\cite{Grunbaum80,Grunbaum85,Grunbaum86,grunbaum1988isonemal}. Multi-layer two-way nonwovens generalize this matrix to a tensor. For fully volumetric nonwovens, each cell of a 3D grid can connect to its neighbors in the positive $x$, $y$, and $z$ directions independently, so each cell's local state is naturally described by an \emph{octal} digit (one bit per axis) rather than a binary one -- directly analogous to, but more general than, the eight local line configurations of a cP-weave cell (Section~\ref{sec:cp-weaves}).

\section{Design: An Algebra of \texorpdfstring{$2^K$}{2\^{}K}-adic and \texorpdfstring{$8^K$}{8\^{}K}-adic Zeros}
\label{sec:padic-design}

Filling a large 3D tensor of local states by hand, or at random, does not scale, and random filling can introduce disconnected ``floating'' components that must be repaired with graph-search post-processing. To design large, symmetric, periodic structures directly, this work develops an algebra over $2^K$-adic zeros -- infinitely repeating $2^K$-adic numbers that all represent zero -- inspired by (1) multiplicative cyclic subgroups of finite fields and (2) classical $p$-adic number systems~\cite{krishnamurthy1975finite}. The key structural fact exploited here is that a repeating decimal (or, more generally, a repeating base-$P$ expansion) is exactly a periodic pattern: just as $0.\overline{3} = 1/3$ and $0.\overline{012}$ and $0.\overline{285714}$ are familiar examples of eventually periodic digit sequences, an infinitely repeating $2^K$-adic (or $8^K$-adic) number is, by construction, a periodic tiling of one-, two-, or three-dimensional space by local configuration labels.

This numeral system supports algebraic operations that act directly on the design:
\begin{itemize}
\item \textbf{Cyclic addition.} Adding a fixed pattern $A$ to a starting pattern $N_x(0)$ repeatedly (modulo the periodicity) generates a whole orbit of related patterns without having to specify each one explicitly.
\item \textbf{Identity and inverse patterns.} The all-zero pattern acts as an additive identity; multiplying a pattern by $-1$ yields its additive inverse, giving a principled way to generate ``complementary'' designs.
\item \textbf{Shift as cyclic multiplication.} Multiplying a $2^K$-adic (or $8^K$-adic) pattern by a power of the base is equivalent to cyclically shifting it -- unifying the ad hoc shift-based notations used earlier for weaves (e.g., the extended $[A/B/s_x/s_y]$ notation) into a single algebraic operation.
\end{itemize}

Because these operations are performed on patterns that are already guaranteed to be periodic (they are, by construction, $2^K$- or $8^K$-adic zeros), every result of every operation is automatically a valid, periodic design -- the algebra never produces a pattern with an infinite, non-repeating fundamental domain. This is the design-time counterpart to the enumeration-time results of Chapter~\ref{cha:volumetric-bravais}: rather than counting how many valid designs exist, the algebra of $P^K$-adic zeros lets a designer \emph{generate} large families of valid, symmetric designs directly, using a handful of algebraic operations (addition, negation, multiplication by a shift) on a compact numeral representation.

\section{Contributions}

In summary, this line of work contributes:
\begin{enumerate}
\item a tensor-based binary/octal numbering scheme that uniquely describes inter-layer and inter-cell connections for two-way, multi-layer, and fully volumetric nonwoven structures;
\item an algebraic framework, based on $2^K$-adic and $8^K$-adic zeros, for constructing tensors that correspond to symmetric nonwoven (and, by the same construction, woven) structures;
\item an identification of the relevant algebraic properties of these numbers (cyclic addition, identity and inverse patterns, shift-as-multiplication), including the observation that a special subset corresponding to Mersenne primes can support cyclic multiplication with additional structure; and
\item a demonstration that the resulting volumetric nonwoven structures can represent previously undocumented infinite regular and semi-regular polyhedra and flexible frame structures~\cite{grunbaum1993infinite,grunbaum2015uniform}.
\end{enumerate}

\section{Relation to the Proposed Framework}

This chapter completes the ``algebraic configuration space'' item identified as a knowledge gap in Chapter~\ref{cha:introduction}: where Chapter~\ref{cha:volumetric-bravais} shows that volumetric weaves \emph{can} be enumerated, this chapter shows that a specific class of them (and their nonwoven relatives) can be \emph{algebraically generated} by design. Chapter~\ref{cha:algebraic-calculus} builds on this algebra to move from single-cell-type, axis-aligned periodic patterns toward a more general, topology-first combinatorial model that also accommodates links, knots, and interlocked (catenated) structures -- moving closer to the fully dimension- and topology-independent design space that is the overall goal of this dissertation (Section~\ref{sec:goal}).
